\documentclass[../../main.tex]{subfiles}% 注意这里的文件路径不能用 ./main.tex ,否则用latexmk编译子文件会报错
\graphicspath{{\subfix{./image/}}} % 指定图片目录,后续可以直接使用图片文件名
% 注意这里的文件路径不能用 ../../image/ ,否则用latexmk编译子文件会报错

% 例如:
% \begin{figure}[H]
% \centering
% \includegraphics[scale=0.3]{图.png}
% \caption{}
% \label{figure:图}
% \end{figure}
% 注意:上述\label{}一定要放在\caption{}之后,否则引用图片序号会只会显示??.

\begin{document}

\section{习题}

\begin{proposition}\label{proposition:Lebesgue外测度的内介值性}
设 \(E \subseteq [a,b]\), \(m^*(E)>0\), \(0 < c < m^*(E)\), 则存在 \(E\) 的有界子集 \(A\), 使得 \(m^*(A)=c\).

特别地,若还有$E$可测,则存在\(E\)的有界可测子集 \(A\), 使得 \(m(A)=c\).
\end{proposition}
\begin{proof}

不妨设$a\geqslant 0$,记 \(f(x)=m^*((-x,x) \cap E)\), \(a \leqslant  x \leqslant  b\), 则 \(f(a)=0\), \(f(b)=m^*(E)\). 对$\forall h>0$,有
\begin{align*}
\left| f\left( x+h \right) -f\left( x \right) \right|&=\left| m^*\left( \left( -x-h,x+h \right) \cap E \right) -m^*\left( \left( -x,x \right) \cap E \right) \right|\\
&=\left| m^*\left[ \left( \left( -x-h,-x \right) \cup \left( x,x+h \right) \right) \cap E \right] \right|\\
&\leqslant m^*\left( -x-h,-x \right) +m^*\left( x,x+h \right) =2h,
\end{align*}
故$f$在$x$右连续.同理，$f$在$x$左连续。从而再由$x$的任意性知，$f\in C\left[ a,b \right]$.
再由$f\left( a \right) =0<a$，$f\left( b \right) =m^*\left( E \right) >a$及连续函数介值定理知，存在$\xi \in \left( a,b \right)$，使得
\begin{align*}
f\left( \xi \right) =m^*\left( \left( -\xi ,\xi \right) \cap E \right) =a.
\end{align*}
取$A=\left( -\xi ,\xi \right) \cap E$即可。显然当$E$可测时,$A$也可测.

\end{proof}

\begin{corollary}
设$(\mathbb{R}^1,\mathscr{L},m)$为Lebesgue测度空间，$E\subset\mathbb{R}^1$。如果$0<a<m(E)$，证明：存在无内点的有界闭集$F\subseteq E$，使得 $m(F)=a$.
\end{corollary}
\begin{remark}
一个集合去掉自身的一个稠密子集就是一个无内点集合.
\end{remark}
\begin{proof}

{\color{blue}证法一:} 令$E_1=E\setminus\mathbb{Q}$，则$E_1$无内点，$m(E_1)=m(E)>a>0$。由\rrefthe{theorem:Lebesgue可测集的充要条件}{theorem:Lebesgue可测集的充要条件-3}知，存在闭集$E_2$，使得 $E_2\subseteq E_1$，且
\begin{align*}
a < m(E_2) < m(E_1) = m(E).
\end{align*}
显然存在$n\in \mathbb{N}$,使得$a+\frac{1}{n}<m(E_2),\forall n>N$.对$\forall n>N$，由\refpro{proposition:Lebesgue外测度的内介值性}同理可知,存在$ r_n\in\mathbb{R}$，使得
\begin{align*}
a < m(E_2\cap[-r_n,r_n]) < a + \frac{1}{n}<m(E_2).
\end{align*}
令$F_n=E_2\cap [-t_n,t_n]$,则$F_n$无内点,有界闭集且
\begin{align*}
F_{n+1}\subseteq F_n,\,\, a<m(F_n)<a+\frac{1}{n},\quad \forall n>N.
\end{align*}
再令$F=\bigcap\limits_{n=1}^{\infty}F_n$，则$F=\bigcap\limits_{n=1}^{\infty}F_n\subseteq E_1\subseteq E$，$F$无内点，有界闭集且
\begin{align*}
a \leqslant m(F) &= m\left(\bigcap\limits_{n=1}^{\infty}F_n\right) \xlongequal{\text{\refthe{corollary:递减可测集列的测度运算}}} \lim_{n\to+\infty} m(F_n) \leqslant \lim_{n\to+\infty} \left(a + \frac{1}{n}\right) = a.
\end{align*}
故$m(F)=a$.

{\color{blue}证法二:} 令$E_1=E\setminus\mathbb{Q}$，则$E_1$无内点，$m(E_1)=m(E)$。由\rrefthe{theorem:Lebesgue可测集的充要条件}{theorem:Lebesgue可测集的充要条件-3}知，存在闭集$E_2$，使得 $E_2\subseteq E_1$，且
\begin{align*}
a < m(E_2) < m(E_1) = m(E).
\end{align*}
显然存在$n\in \mathbb{N}$,使得$a+\frac{1}{n}<m(E_2),\forall n>N$.对$\forall n>N$，由\refpro{proposition:Lebesgue外测度的内介值性}同理可知,存在$ r_n\in\mathbb{R}$，使得
\begin{align*}
a < m(E_2\cap[-r_n,r_n]) < a + \frac{1}{n}<m(E_2).
\end{align*}
令$r=\lim_{n\to+\infty} r_n$，则
\begin{align*}
a \leqslant m(E_2\cap[-r,r]) = \lim_{n\to+\infty} m(E_2\cap[-r_n,r_n]) \leqslant \lim_{n\to+\infty} \left(a + \frac{1}{n}\right) = a,
\end{align*}
于是$m(E_2\cap[-r,r]) = a.$,故$F=E_2\cap[-r,r]$为无内点的有界闭集，且$F\subseteq E$，$m(F)=a$.

{\color{blue}证法三:} 令$E_1=E\setminus\mathbb{Q}$，则$E_1$无内点，$m(E_1)=m(E)$。由\rrefthe{theorem:Lebesgue可测集的充要条件}{theorem:Lebesgue可测集的充要条件-3}知，存在闭集$F_1\subseteq E_1\subseteq E$，使得 $m(F_1)>a$。再作
$$
f:\mathbb{R}^+=[0,+\infty)\to\mathbb{R},\,\,x\mapsto f(x)=m(F_1\cap[-x,x]).
$$
由
\begin{align*}
|f(x)-f(y)| \leqslant 2|x-y|
\end{align*}
立知，$f\in C(\mathbb{R}^+)$，且$f(0)=0$，$\lim_{x\to+\infty} f(x)=m(F_1)$。根据连续函数的介值定理，$\exists x_0\in\mathbb{R}^+$，使得 $f(x_0)=m(F_1\cap[-x_0,x_0])=a$。于是，$F=F_1\cap[-x_0,x_0]$为所求的无内点的有界闭集，且$F\subseteq E$，$m(F)=a$.

{\color{blue}证法四:} 因为$0<a<m(E)$，故存在有界闭集$F_1\subseteq E$，使得 $m(F_1)=b>a$。令
\begin{align*}
G = \bigcup\limits_{n=1}^{\infty} \left(r_n - \frac{\alpha}{n^2}, r_n + \frac{\alpha}{n^2}\right),
\end{align*}
其中$\{r_n\}$为$\mathbb{Q}$的一个排列。取$0<\alpha<\frac{3(b-a)}{\pi^2}$，$F_1\cap G^c$为无内点的闭集，且
\begin{align*}
m(F_1\cap G^c) = m(F_1) - m(F_1\cap G) 
\geqslant b - m(G) = b - \alpha \cdot \sum\limits_{n=1}^{\infty} \frac{2}{n^2} 
> b - \frac{3(b-a)}{\pi^2} \cdot \frac{\pi^2}{3} = a.
\end{align*}
根据\refpro{proposition:Lebesgue外测度的内介值性},存在有界闭集$K=[-t,t]$，使得
\begin{align*}
m([-t,t]\cap F_1\cap G^c) = a.
\end{align*}
于是，$F=K\cap F_1\cap G^c$为无内点的有界闭集，且$F\subseteq F_1\subseteq E$，$m(F)=a$.

\end{proof}

\begin{proposition}
对$\forall E\subset\mathbb{R}^n$，证明：存在$F_\sigma$集$K$和$G_\delta$集$H$，使得$K\subseteq E\subseteq H$，且$m_*(E)=m(K)$，$m^*(E)=m(H)$。
\end{proposition}
\begin{proof}

设$E\in\mathbb{R}^n=\mathscr{K}(\mathscr{R}_0)$，其中$\mathscr{R}_0$为半开半闭长方体的有限并的全体。根据$m^*$的定义，$\exists E_k\in\mathscr{R}_0$（$k\in\mathbb{N}$），使得$E\subseteq\bigcup\limits_{k=1}^{\infty}E_k$，$\sum\limits_{k=1}^{\infty}m(E_k)<m^*(E)+\frac{1}{l}$。显然，必有开集$H_l$，使得$E\subseteq\bigcup\limits_{k=1}^{\infty}E_k\subseteq H_l$，且$m(H_l)<m^*(E)+\frac{1}{l}$（只需将每个半开半闭的长方体稍扩大一点得到$H_l$）。

同样，根据$m_*$的定义，$\exists F_l\in\mathscr{R}_0$，使得$F_l\subseteq E$，$m(F_l)>m_*(E)-\frac{1}{l}$。显然，必有闭集$K_l$，使得$K_l\subseteq F_l\subseteq E$，$m(K_l)>m_*(E)-\frac{1}{l}$（只需将每个半开半闭的长方体稍缩小一点得$K_l$）。
\begin{gather*}
m_*(E)-\frac{1}{l} \leqslant m(K_l)\leqslant m\left(\bigcup\limits_{l=1}^{\infty}K_l\right)\leqslant m_*(E), \\
m^*(E) \leqslant m\left(\bigcap\limits_{l=1}^{\infty}H_l\right)\leqslant m(H_l)<m^*(E)+\frac{1}{l},
\end{gather*}
所以，当$l\to+\infty$时得到
\begin{align*}
m\left(\bigcup\limits_{l=1}^{\infty}K_l\right) = m_*(E),\quad m\left(\bigcap\limits_{l=1}^{\infty}H_l\right) = m^*(E).
\end{align*}
我们令$K=\bigcup\limits_{l=1}^{\infty}K_l$，$H=\bigcap\limits_{l=1}^{\infty}H_l$，则$K$，$H$分别为题中要求的$F_\sigma$集和$G_\delta$集。

\end{proof}

\begin{example}
设$(\mathbb{R}^1,\mathscr{L},m)$为Lebesgue测度空间。
\begin{enumerate}[(1)]
\item 作出一个由$\mathbb{R}^1$中某些无理数构成的闭集$F$，使得$m(F)>0$.
\item 设有理数集$\mathbb{Q}=\{r_n\mid n\in\mathbb{N}\}$，令
\[
G = \bigcup\limits_{n=1}^{\infty}\left(r_n - \frac{1}{n^2}, r_n + \frac{1}{n^2}\right).
\]
证明：任一闭集$F\subset\mathbb{R}^1$，有$m(G\Delta F)>0$.
\end{enumerate}
\end{example}
\begin{solution}
\begin{enumerate}[(1)]
\item {\color{blue}解法一:} 记$\mathbb{Q}=\{r_1,r_2,\cdots\}$，令$G_n=\left(r_n-\frac{1}{2^n},r_n+\frac{1}{2^n}\right)$，则$G=\bigcup\limits_{n=1}^{\infty}G_n=\bigcup\limits_{n=1}^{\infty}\left(r_n-\frac{1}{2^n},r_n+\frac{1}{2^n}\right)$为开集，且$m(G)\leqslant\sum\limits_{n=1}^{\infty}\frac{2}{2^n}=2$。令$F=G^c$，则$F$为闭集，且$m(F)=m(\mathbb{R}^1)-m(G)\geqslant(+\infty)-2=+\infty>0$.

{\color{blue}解法二:} 设$E=[0,1]\setminus\mathbb{Q}$，则$m(E)=1$。由\rrefthe{theorem:Lebesgue可测集的充要条件}{theorem:Lebesgue可测集的充要条件-3}知,存在$F\subset E$为闭集，且$m(E\setminus F)<\frac{1}{2}$。因此$F$全由无理数组成,且
\begin{align*}
m(F) = m(E)-m(E\setminus F)\geqslant 1-\frac{1}{2}=\frac{1}{2}>0.
\end{align*}
\end{enumerate}
\begin{enumerate}[(1)]
\setcounter{enumi}{1}
\item {\color{blue}证法一:}设$F\subset\mathbb{R}^1$为闭集，若$F\supseteq G$，则闭集$F\supseteq\overline{G}=\mathbb{R}^1$，$F=\mathbb{R}^1$。自然有
\begin{align*}
m(G\Delta F) &= m(G\setminus F)+m(F\setminus G) = m(\varnothing)+m(\mathbb{R}^1\setminus G) \\
&= 0+m(\mathbb{R}^1)-m(G) = (+\infty)-\sum\limits_{n=1}^{\infty}\frac{2}{n^2}=+\infty>0.
\end{align*}
若$F\nsupseteq G$，则$\exists x_0\in G\setminus F$，由$G$为开集，$F$为闭集知，$G\setminus F=G\cap F^c$为开集，故$\exists\delta>0$，使得$B(x_0;\delta)\subseteq G\setminus F$。因此
\begin{align*}
m(G\Delta F)\geqslant m(G\setminus F)\geqslant m(B(x_0;\delta))\geqslant 2\delta>0.
\end{align*}

{\color{blue}证法二:} (反证)假设存在闭集$F$，使得$m(G\setminus F)+m(F\setminus G)=m(G\Delta F)=0$，则
\[
m(G\setminus F)=0,\ m(F\setminus G)=0,\ m(F)=m(F\cap G)=m(G).
\]
易证$G\subseteq F$（事实上，$\forall x_0\in G$，若$x_0\notin F$，则由$G$为开集，而$F$为闭集知，$\exists\delta_0>0$，使得$B(x_0;\delta_0)\subseteq G$且$B(x_0;\delta_0)\cap F=\varnothing$。从而
\[
0 = m(G\setminus F)\geqslant m(B(x_0;\delta_0))=2\delta_0>0,
\]
矛盾。因此，$G\subseteq F$。于是，$\mathbb{R}^1=\overline{G}\subseteq\overline{F}=F$，从而
\begin{gather*}
m(F)\geqslant m(\mathbb{R}^1)=+\infty, \\
+\infty>\frac{\pi^2}{3}=\sum_{n=1}^{\infty}{\frac{2}{n^2}}\geqslant m(G)=m(F)\geqslant+\infty,
\end{gather*}
矛盾。

{\color{blue}证法三:}(反证)假设存在闭集$F$，使得$m(G\setminus F)+m(F\setminus G)=m(G\Delta F)=0$，则
\[
m(G\setminus F)=0,\quad m(F\setminus G)=0.
\]
于是，开集$G\cap F^c=G\setminus F=\varnothing$(否则由$G\setminus F$是开集知,存在开球$B(x_0,\delta)\subseteq G\setminus F,$这与$m(G\setminus F)=0$矛盾)。从而，$G\subseteq F$，$\mathbb{R}^1=\overline{G}\subseteq\overline{F}=F$，$F=\mathbb{R}^1$。此时
\begin{align*}
0 = m(F\setminus G)=m(\mathbb{R}^1\cap G^c)=m(G^c).
\end{align*}
这与
\begin{align*}
m(G^c)=m(\mathbb{R}^1)-m(G)=(+\infty)-\sum\limits_{n=1}^{\infty}\frac{2}{n^2}=(+\infty)-\frac{\pi^2}{3}=+\infty
\end{align*}
相矛盾。
\end{enumerate}

\end{solution}

\begin{example}
设$(\mathbb{R}^n,\mathscr{L},m)$为Lebesgue测度空间。
\begin{enumerate}[(1)]
\item 设$E\subset\mathbb{R}^n$，且$m^*(E)<+\infty$。如果
\[
m^*(E) = \sup\{m(F) \mid F\subset E\text{ 为有界闭集}\},
\]
证明：$E\in\mathscr{L}$。
如果条件“$m^*(E)<+\infty$”删去，结论不一定正确。
\item 设$E\in\mathscr{L}$且$m(E)>0$。证明：存在$x\in E$，使得对$\forall\delta>0$，有
\[
m(E\cap B(x;\delta))>0.
\]
\item 设$\{B_k\}$为$\mathbb{R}^n$中递减可测集列，$m^*(A)<+\infty$。令
\[
E_k = A\cap B_k,\ k\in\mathbb{N},\quad E = \bigcap_{k=1}^{\infty}E_k.
\]
证明：$\lim_{k\to+\infty}m^*(E_k)=m^*(E)$。
\end{enumerate}
\end{example}
\begin{remark}
第(1)问注意
\begin{align*}
m^*(E) &\xlongequal{\text{题设}} \sup\{m(F) \mid F\subset E\text{ 为有界闭集}\} \\
&\leqslant \sup\{m(\widetilde{F}) \mid \widetilde{F}\subset E\text{ 中的闭集}\} \\
&= \sup\{m(\widetilde{F}\cap[-k,k]) \mid \widetilde{F}\subset E\text{ 中的闭集},k\in\mathbb{N}\} \\
&\leqslant \sup\{m(F) \mid F\subset E\text{ 为有界闭集}\} = m^*(E), \\
m^*(E) &= \sup\{m(F) \mid F\subset E\text{ 为有界闭集}\} = \sup\{m(\widetilde{F}) \mid \widetilde{F}\subset E\text{ 为闭集}\}.
\end{align*}
\end{remark}
\begin{proof}
\begin{enumerate}[(1)]
\item {\color{blue}证法一:} 因为
\[
m^*(E) = \sup\{m(F) \mid F\subset E\text{ 为有界闭集}\},
\]
所以存在有界闭集$F_k\subset E$，使得$m(F_k)>m^*(E)-\frac{1}{k}$。于是
\begin{align*}
m^*(E)-\frac{1}{k} \leqslant m(F_k) \leqslant m^*\left(\bigcup_{i=1}^{\infty}F_i\right) \leqslant m^*(E).
\end{align*}
令$k\to+\infty$，得到
\begin{align*}
m^*\left(\bigcup_{i=1}^{\infty}F_i\right) = m^*(E)<+\infty.
\end{align*}
再根据$m^*$定义知，存在开集$G_k\supset E$，使得$m(G_k)<m^*(E)+\frac{1}{k}$，
\begin{align*}
m^*(E) \leqslant m^*\left(\bigcap_{i=1}^{\infty}G_i\right) \leqslant m(G_k) < m^*(E)+\frac{1}{k}.
\end{align*}
令$k\to+\infty$，得到
\begin{align*}
m^*\left(\bigcap_{i=1}^{\infty}G_i\right) = m^*(E).
\end{align*}
由此，有$\bigcup_{i=1}^{\infty}F_i \subset E \subset \bigcap_{i=1}^{\infty}G_i$,于是由\rrefthe{theorem:测度的基本性质}{theorem:测度的基本性质-2}得
\begin{gather*}
m\left(\bigcap_{i=1}^{\infty}G_i - \bigcup_{i=1}^{\infty}F_i\right) = m\left(\bigcap_{i=1}^{\infty}G_i\right)-m\left(\bigcup_{i=1}^{\infty}F_i\right) = m^*(E)-m^*(E)=0, \\
m^*\left( \bigcap_{i=1}^{\infty}{G_i}-E \right) \leqslant m\left( \bigcap_{i=1}^{\infty}{G_i}-\bigcup_{i=1}^{\infty}{F_i} \right) =0.
\end{gather*}
从而$\bigcap_{i=1}^{\infty}{G_i}-\bigcup_{i=1}^{\infty}{F_i},\,\,-E\in \mathscr{L} .$故
\begin{align*}
E = \bigcap_{i=1}^{\infty}G_i - \left(\bigcap_{i=1}^{\infty}G_i - E\right)\in\mathscr{L}.
\end{align*}

{\color{blue}证法二:} $\forall T\subset\mathbb{R}^n$，对有界闭集$F\subset E$，有$F^c\supset E^c$且
\begin{align}
m^*(T) = m^*(T\cap F)+m^*(T\cap F^c) \geqslant m^*(T\cap F)+m^*(T\cap E^c). \label{eq::20fjsuf09j3f30jjwfawf23r32t--}
\end{align}
由\rrefthe{theorem:测度的基本性质}{theorem:测度的基本性质-2}得
\begin{align*}
0\leqslant m^*(T\cap E)-m^*(T\cap F) \leqslant m^*(T\cap (E-F)) \leqslant m^*(E-F) 
 = m^*(E)-m(F).
\end{align*}
因为$m^*(E)<+\infty$，所以
\begin{align*}
m^*(T\cap E)-m^*(E)+m(F) \leqslant m^*(T\cap F) \leqslant m^*(T\cap E).
\end{align*}
就有界闭集$F\subset E$取$\sup$可得
\begin{align*}
m^*(T\cap E) = m^*(T\cap E)-m^*(E)+m^*(E) \leqslant \sup_F m^*(T\cap F) \leqslant m^*(T\cap E).
\end{align*}
因此
\begin{align*}
m^*(T\cap E) = \sup_F m^*(T\cap F).
\end{align*}
再在\eqref{eq::20fjsuf09j3f30jjwfawf23r32t--}式取$\sup$得
\begin{align*}
m^*(T) \geqslant \sup_F m^*(T\cap F)+m^*(T\cap E^c) = m^*(T\cap E)+m^*(T\cap E^c).
\end{align*}
另一方面，由$m^*$的次可加性，显然有
\begin{align*}
m^*(T) \leqslant m^*(T\cap E)+m^*(T\cap E^c).
\end{align*}
综上得到
\begin{align*}
m^*(T) = m^*(T\cap E)+m^*(T\cap E^c).
\end{align*}
这就证明了$E$为Lebesgue可测集。

{\color{blue}证法三:} 因为
\[
m^*(E) = \sup\{m(F) \mid F\subset E\text{ 为有界闭集}\}<+\infty,
\]
所以，对$\forall\varepsilon>0$，$\exists F\subset E$为有界闭集，使得
\[
m^*(E)-m(F)<\varepsilon.
\]
而由闭集$F$为Lebesgue可测集，故
\begin{gather*}
m^*(E) = m^*(E\cap F)+m^*(E\cap F^c) = m^*(F)+m^*(E\cap F^c), \\
m^*(E\cap F^c) = m^*(E)-m^*(F) < \varepsilon.
\end{gather*}
对$\forall T\subset\mathbb{R}^n$，有
\begin{align*}
&m^*(T) \leqslant m^*(T\cap E)+m^*(T\cap E^c) \\
&= m^*((T\cap F)\cup(T\cap E\cap F^c))+m^*(T\cap E^c) \\
&\leqslant m^*(T\cap F)+m^*(T\cap E\cap F^c)+m^*(T\cap E^c) \\
&\leqslant m^*(T\cap F)+m^*(E\cap F^c)+m^*(T\cap F^c) < m^*(T)+\varepsilon.
\end{align*}
令$\varepsilon\to 0^+$，得到
\begin{align*}
m^*(T) = m^*(T\cap E)+m^*(T\cap E^c).
\end{align*}
这就证明了$E$为Lebesgue可测集。

如果条件“$m^*(E)<+\infty$”删去，结论不一定正确。

反例：设$E=[0,+\infty)\cup W$，$W$为$(-\infty,0)$中的不可测集，$F_k=[0,k]$为有界闭集，$m([0,k])=k$，
\[
m^*(E)=+\infty = \sup\{m(F) \mid F\subset E\text{ 为有界闭集}\},
\]
但$E$为Lebesgue不可测集。

\item (反证)反设$\forall x\in E$，$\exists\delta(x)>0$，使得
\[
m(E\cap B(x;\delta(x)))=0.
\]
显然，$\{B(x;\delta(x))\mid x\in E\}$为$E$的一个开覆盖。根据\hyperref[theorem:Lindelöf定理]{Lindelöf定理}，它有可数的子覆盖
\[
\{B(x_i;\delta(x_i)) \mid i\in\mathbb{N},x_i\in E\}.
\]
因此
\begin{align*}
0 < m(E) &= m\left(E\cap\left(\bigcup_{i=1}^{\infty}B(x_i;\delta(x_i))\right)\right) \\
&\leqslant \sum_{i=1}^{\infty}m(E\cap B(x_i;\delta(x_i))) = \sum_{i=1}^{\infty}0 = 0,
\end{align*}
矛盾.

\item {\color{blue}证法一:} 由于$B_k$可测，故
\[
m^*(A) = m^*(A\cap B_k)+m^*(A-B_k).
\]
由于$\bigcap_{k=1}^{\infty}B_k$仍为可测集，所以
\begin{align}\label{eq:::--2fhr89e9f}
m^*(A) = m^*\left(A\cap\left(\bigcap_{k=1}^{\infty}B_k\right)\right)+m^*\left(A-\bigcap_{k=1}^{\infty}B_k\right).
\end{align}
因为$B_k$递减，故$A-B_k$递增，根据\refcor{corollary:递增集列外测度和极限可交换},有
\begin{align*}
\lim_{k\to+\infty}m^*(A-B_k) = m^*\left(\lim_{k\to+\infty}(A-B_k)\right) = m^*\left(\bigcup_{k=1}^{\infty}(A-B_k)\right) \xlongequal{\text{de Morgan公式}} m^*\left(A-\bigcap_{k=1}^{\infty}B_k\right).
\end{align*}
由此推得（注意$m^*(A)<+\infty$）
\begin{align*}
\lim_{k\to+\infty}m^*(E_k) &= \lim_{k\to+\infty}m^*(A\cap B_k) = \lim_{k\to+\infty}\left[m^*(A)-m^*(A-B_k)\right] \\
&= m^*(A)-m^*\left(\bigcup_{k=1}^{\infty}(A-B_k)\right) \xlongequal{\text{de Morgan公式}} m^*(A)-m^*\left(A-\bigcap_{k=1}^{\infty}B_k\right) \\
&\xlongequal{\text{见\eqref{eq:::--2fhr89e9f}式}} m^*\left(A\cap\left(\bigcap_{k=1}^{\infty}B_k\right)\right) = m^*\left(\bigcap_{k=1}^{\infty}(A\cap B_k)\right) = m^*\left(\bigcap_{k=1}^{\infty}E_k\right) = m^*(E).
\end{align*}

{\color{blue}证法二:} 因为$B_k$递减可测，故$B_k^c = A-B_k = A\cap B_k^c$递增，且$\bigcap_{k=1}^{\infty}B_k$仍可测。于是
\begin{align*}
&m^*(A)-\lim_{k\to+\infty}m^*(E_k) = m^*(A)-\lim_{k\to+\infty}m^*(A\cap B_k) \\
&= \lim_{k\to+\infty}\left[m^*(A)-m^*(A\cap B_k)\right] = \lim_{k\to+\infty}m^*(A\cap B_k^c) \\
&\xlongequal{\text{\refcor{corollary:递增集列外测度和极限可交换}}} m^*\left(\lim_{k\to+\infty}(A\cap B_k^c)\right) = m^*\left(\bigcup_{k=1}^{\infty}(A\cap B_k^c)\right) \\
&= m^*\left(A\cap\left(\bigcup_{k=1}^{\infty}B_k^c\right)\right) \xlongequal{\text{de Morgan公式}} m^*\left(A\cap\left(\bigcap_{k=1}^{\infty}B_k\right)^c\right) \\
&\xlongequal{\bigcap_{k=1}^{\infty}B_k\text{ 可测}} m^*(A)-m^*\left(A\cap\left(\bigcap_{k=1}^{\infty}B_k\right)\right) = m^*(A)-m^*\left(\bigcap_{k=1}^{\infty}(A\cap B_k)\right) \\
&= m^*(A)-m^*\left(\lim_{k\to+\infty}(A\cap B_k)\right) = m^*(A)-m^*\left(\lim_{k\to+\infty}E_k\right) \\
&= m^*(A)-m^*(E),
\end{align*}
注意到$m^*(A)<+\infty$，因此有
\begin{align*}
\lim_{k\to+\infty}m^*(E_k) = m^*(E).
\end{align*}
\end{enumerate}

\end{proof}






















\end{document}